\clearpage
\section{How the comparison is made fair}\label{sec:design}
\subsection{Specify the experiment before looking for a winner}
A machine-learning result combines choices about data, representation, objective,
optimization, model selection, and evaluation. If those choices are changed
after seeing a preferred method's outcomes, apparent improvement can reflect
the selection process rather than a reliable advantage.

The v3 protocol fixes the hypothesis, action splits, gallery, recipes, duration,
tuning grid, corruption conditions, analysis, and reporting rules before
selection and final-test access. This is a \emph{prespecified} comparison. A
hash-bound configuration makes the exact specification checkable; it does not
replace a scientific explanation of why the choices are appropriate.

\begin{longtable}{@{}P{0.28\linewidth}P{0.32\linewidth}P{0.34\linewidth}@{}}
\toprule Controlled choice & What stays the same & Alternative explanation reduced\\\midrule\endhead
Representation & Encoder, checkpoint, preprocessing, head structure & A different
backbone or input pipeline caused the difference.\\\addlinespace
Optimization exposure & Twenty epochs, physical $8\times4$ batches & One recipe
received more updates or different within-batch information.\\\addlinespace
Initialization and order & Exact initial weights and ordered batches within seed
& A favorable starting state or example order caused the difference.\\\addlinespace
Tuning budget & Three rates, three seeds per rate, same ranking rule & More search
opportunities gave one method an advantage.\\\addlinespace
Evaluation input & Same gallery, queries, and per-example damage & Methods were
evaluated on different problems.\\\addlinespace
Claim rule & One primary effect and fixed interval construction & A favorable
metric, subset, or threshold was chosen afterward.\\\bottomrule
\end{longtable}

These controls improve \textbf{internal validity}: whether the observed
difference answers the declared comparison. \textbf{External validity} asks
whether it transfers to other datasets, people, actions, cameras, or damage
mechanisms. Internal controls alone cannot establish that broader claim.

\subsection{Pairing is more than using the same seed}
Let $A_s$ and $B_s$ denote outcomes from two methods under a paired condition.
For their difference,
\[
 \operatorname{Var}(A_s-B_s)
 =\operatorname{Var}(A_s)+\operatorname{Var}(B_s)
  -2\operatorname{Cov}(A_s,B_s).
\]
Shared difficulty can induce positive covariance, making the difference less
variable than an unpaired comparison. Pairing is useful even when that variance
reduction is not guaranteed: it prevents known input differences from becoming
uncontrolled method differences.

A seed initializes a pseudorandom generator. It does not prove two execution
paths consumed random values in the same order or used the same arithmetic.
The project therefore stores and verifies exact starting-state and batch-plan
identities. Seeds 7, 17, and 29 give three paired training realizations.
They produce three separate heads per method, \emph{not an ensemble}.

An epoch is defined as $\lceil N/32\rceil$ balanced batches. For candidate
training, $\lceil76{,}013/32\rceil=2{,}376$ updates per epoch and 47,520 over
twenty epochs. Class and example queues are shuffled at the epoch start and
cycled within it. This is a balanced update budget, not a guarantee that every
row is visited exactly once per epoch.

Gradient accumulation cannot replace the physical batch: relationships between
examples in different microbatches never exist in either microbatch's similarity
matrix. Adding their gradients cannot retroactively construct the missing
neighbors, positives, or softmax alternatives.

\subsection{Pilots, model selection, and final training are separate}
\textbf{Engineering pilots} answer whether the implementation can complete the
specified training and produce valid numerical artifacts. There is one pilot
per recipe, at seed 7 and learning rate $3\times10^{-4}$, for twenty epochs.
Checkpoints at epochs 5, 10, 15, and 20 plus epoch-zero diagnostics expose the
training behavior. Their accuracy does not select a learning rate or a winning
method. A pilot that happens to share a candidate's settings is still ineligible
for selection.

\textbf{Model selection} chooses the learning rate separately for each recipe:
\[
 3\text{ recipes}\times3\text{ rates}\times3\text{ seeds}=27\text{ fresh runs},
 \qquad\eta\in\{10^{-4},3\times10^{-4},10^{-3}\}.
\]
Rank candidates using unrounded mean epoch-20 clean development top-1 accuracy
over all three seeds. Resolve exact ties with mean MRR, then lower learning rate.
The best checkpoint is not selected across epochs. Corrupted development scores
are not tuning criteria. If an endpoint learning rate wins, the grid is not
expanded after the fact.

\textbf{Final training} starts nine fresh heads: three recipes at their selected
rates, each with three paired seeds, using all 95,001 auxiliary examples. All nine
must finish and be locked before novel evaluation. A final run is not a pilot
continued for longer, a selected checkpoint reused unchanged, or a seed chosen
because it looked favorable.

\begin{concept}{What a failed run means}
A candidate needs all three valid seeds; averaging only its two successful
seeds changes the comparison. Preserve failures and their causes. An allowed
infrastructure retry repeats identical inputs in a new attempt; it is not a
chance to draw a better seed. Numerical failure and poor accuracy are different:
poor but valid accuracy is a scientific outcome, while nonfinite computation
invalidates the affected evidence.
\end{concept}

The fixed twenty-epoch duration avoids choosing a favorable stopping point.
It does not claim that every method has converged or received its ideal budget.
The fair statement is equal declared training and selection opportunities,
with cost and limitations reported.

\subsection{Why keep the final test closed?}
Repeatedly looking at test results can influence human choices even if no
test gradient is computed. Picking margins, a split, a corruption, or a headline
after observing novel outcomes would use the test for development.

The study therefore progresses through separate authorizations: auxiliary
engineering, selection, final-head locking, and only then the single novel
evaluation campaign. A dataset-level opening record spans historical and current
versions so a new output folder cannot silently reset the holdout's history.
After opening, scientific training and selection for this study/source remain
closed. An error discovered afterward must be preserved and reported; a corrected
analysis cannot silently replace the confirmatory result.

There is a historical packaging nuance. The original aggregate loader
deserialized all annotations before selecting auxiliary records. The evidence
supports no novel retrieval preprocessing, embedding, visualization, scoring,
or tuning access; it does not support the stronger claim that no novel bytes
were ever deserialized. The newer auxiliary package separates authorized inputs
from novel processing. State the actual access boundary, not a cleaner-sounding
history.

\source{\src{plan/research-plan.v3.md}{Plan, Sections 4 and 9};
\src{src/pose_embed/training/sampler.py}{Balanced sampler};
\src{docs/protocol/protocol-v3.md}{Fixed recipes, selection, and staged authorization}.}

\section{Designing a controlled disturbance}\label{sec:corruption}
\subsection{Where the intervention occurs}
Write preprocessing as $P$, frozen encoding and pooling as $F$, and the learned
head as $H_m$. Clean retrieval uses $H_m(F(P(x)))$. A damaged query uses
\[
 z_{m,c}=H_m\bigl(F(T_c(P(x)))\bigr),
\]
where $T_c$ changes the prepared pose tensor. The gallery remains clean.
This location is part of the scientific question: the study perturbs pose
observations, not rendered pixels, model weights, training labels, or cached
embedding coordinates.

Each family starts from its own clean tensor. Damage families are not combined.
After damage, there is no clipping, recentering, normalization, interpolation,
or frame deletion that could alter the declared intervention.

\begin{tabularx}{\linewidth}{@{}P{0.20\linewidth}P{0.24\linewidth}Y@{}}
\toprule Family & Three strengths & What changes and what stays fixed\\\midrule
Coordinate jitter & 1\%, 2.5\%, 5\% of torso scale & Gaussian offsets to observed
$x,y$ coordinates; confidence and missing observations unchanged.\\\addlinespace
Joint masking & 3, 6, 8 joints & Zero complete trajectories, including confidence,
using nested prefixes of one anatomically grouped ordering.\\\addlinespace
Frame masking & 10, 25, 40 of 100 frames & Zero a contiguous nested interval across
people, joints, and channels; retain all 100 time slots.\\\bottomrule
\end{tabularx}

\subsection{Scale jitter to the body, not arbitrary pixels}
For observed joint coordinates, the intervention has the form
\[
 x'_{p,t,j,:2}=x_{p,t,j,:2}+\sigma\,r(x)\,\eta_{p,t,j},
 \qquad\eta_{p,t,j}\sim\mathcal N(0,I_2).
\]
$\sigma$ is a fixed severity and $r(x)$ is the clean sample's lower-median
positive finite torso scale. Thus the noise standard deviation is expressed
relative to body size. A sample without usable torso geometry uses one fixed
fallback: the lower median of valid sample-level scales from the eighty-action
training set. It is measured and frozen before the relevant runs, not estimated
from final-test poses or tuned to produce convenient damage.

The same Gaussian field is reused at all three severities, multiplied by the
corresponding scale. That keeps the disturbance direction fixed while changing
its magnitude. Missing or padded observations remain missing rather than being
turned into plausible new joints.

\subsection{Nested severity and shared randomness}
Per-example randomness is determined by the sample identity and corruption
family, not by method, training seed, or severity. Joint masks use nested
prefixes; frame intervals are coupled so larger intervals contain smaller ones.
Therefore, when severity increases, it strengthens the same intervention.
It does not accidentally replace it with an unrelated easier random pattern.

The same damaged tensor is evaluated for every recipe and seed. A comparison
then measures different responses to the same disturbance. It does not compare
one model's fortunate noise draw with another model's unfortunate draw.

\subsection{Validate the intervention before interpreting performance}
A declared ``40-frame mask'' should zero the intended forty slots without
shortening the sequence. Jitter should preserve confidence and scale with the
declared body measure. The geometry panel checks interpretable examples and edge
cases; full-query damage summaries establish what happened across the evaluation
population, including fallback use and zero-damage cases.

That evidence is a manipulation check: did $T_c$ do what the experiment says?
It does not establish that contextual training handles the damage better.
The October development damage audit passed its recorded execution/damage
checks, but later runtime-dependent evidence still requires regeneration.

These are controlled tensor stress tests. Real occlusion may cause a pose
estimator to hallucinate confident joints rather than set them to zero; real
frame loss may have another temporal pattern. MotionBERT was already pretrained
with masking/noise. The present experiment studies the incremental effect of
head-training recipes on that shared encoder.

\studycheck{Why exclude severity from the random seed? Why measure fallback
scale only from training poses? Why is masking input coordinates not equivalent
to simulating a camera occlusion? Answers: Questions 8--9.}
\source{\src{docs/protocol/protocol-v3.md}{Corruption conventions};
\src{src/pose_embed/corruptions/pose.py}{Corruption operators};
\src{src/pose_embed/corruptions/fallback.py}{Training-only fallback};
\src{docs/protocol/week-5-execution.md}{Recorded development damage audit}.}

\section{From retrieval ranks to a defensible conclusion}\label{sec:analysis}
\subsection{Compute ranks first, then metrics}
Let $r_q$ be the rank of the one correct gallery exemplar for query $q$.
With a fixed gallery order to break exact cosine ties,
\[
 \mathrm{Top1}=\frac1Q\sum_q\ind[r_q=1],\qquad
 \mathrm{R@5}=\frac1Q\sum_q\ind[r_q\le5],\qquad
 \mathrm{MRR}=\frac1Q\sum_q\frac1{r_q}.
\]
For invented ranks $(1,2,7,1)$, Top-1 is $2/4=0.5$, R@5 is $3/4=0.75$, and
MRR is $(1+1/2+1/7+1)/4\approx0.661$. Top-1 asks whether the immediate answer
is correct. R@5 asks whether the correct answer is in a short list. MRR rewards
moving the correct answer closer to the top even when it is not first.

With exactly one relevant gallery item, average precision per query is $1/r_q$,
so mAP equals MRR. Reporting both would not add independent evidence.

The complete final report has
\[
 3\text{ recipes}\times3\text{ seeds}\times10\text{ conditions}
 \times2\text{ query definitions}=180\text{ cells}.
\]
The ten conditions are clean plus nine corruptions. These are 180 reporting
cells, not 180 independent experiments: the same heads, examples, and damage
draws recur across them. The two query sets also overlap heavily.

\subsection{Derive the primary effect at the query level}
Let $a_{m,s,c,q}\in\{0,1\}$ indicate a correct top-1 prediction. For each
primary query, compute
\[
 e_q=\frac1{27}\sum_{s,c}
 \left[(a_{C,s,0,q}-a_{C,s,c,q})
       -(a_{B,s,0,q}-a_{B,s,c,q})\right],
\]
where $C$ is contextual and $B$ is standard contrastive. Then
\[
 \Delta=\frac1Q\sum_q e_q.
\]
This equals the earlier difference of mean accuracy drops because sums can be
reordered. A query that becomes correct after corruption contributes a negative
drop; the definition does not clip improvements to zero.

Queries have equal weight within each cell. The three seeds and nine corruption
conditions have equal weight in the declared mixture. This does not imply the
nine conditions are equally common in the world. It defines the specific stress
mixture whose average effect is being tested.

An algebraic check connects the headline to absolute performance:
\[
 \Delta=(A_{C,\mathrm{clean}}-A_{B,\mathrm{clean}})
        -(\overline A_{C,\mathrm{corrupt}}-\overline A_{B,\mathrm{corrupt}}).
\]
Always report all three quantities. An accuracy difference of $0.03$ is
\emph{three percentage points}, not a three-percent relative increase. It
corresponds descriptively to thirty additional correct matches per thousand
queries under the specified mixture.

\subsection{Why ordinary row resampling would be misleading}
The observed $\Delta$ is a number for this fixed collection of trained heads
and examples. To characterize sampling variation within that setting, the
protocol uses a bootstrap: repeatedly resample the observed units with
replacement and recompute the statistic.

The resampling unit must reflect dependence. If several cameras recorded the
same performance, resampling each view independently treats correlated evidence
as if it came from separate events. The project instead samples entire
performance groups, preserving all available camera views together.

The declared procedure is:
\begin{enumerate}
\item Within each of the twenty fixed actions, identify the performance groups.
\item Sample that action's original number of groups with replacement. A group
selected twice contributes all its queries twice.
\item Use the same sampled groups across all methods, seeds, and conditions.
\item Sum their per-query effects and divide by the total sampled query count.
\item Repeat 10,000 times with the fixed generator seed, then take the 2.5th and
97.5th percentiles using the prescribed linear quantile convention.
\end{enumerate}
Actions and training seeds are \emph{not} resampled. This is not a bootstrap
over all possible tasks or all possible training runs.

\subsection{Unequal camera counts change the denominator}
Suppose one sampled performance group contains one query with effect $+1$ and
another contains three queries each with effect $-1$. The query-weighted mean is
\[
 \frac{1-1-1-1}{4}=-0.5.
\]
The mean of the two group means would be $(1+(-1))/2=0$, which answers a
different question. Groups are resampled together to respect dependence, but
queries retain their declared weighting. Resampling is not permission to change
the estimand, meaning the exact quantity we intend to estimate.

Because groups can have different numbers of views, each bootstrap draw can
contain a different query count. The denominator must be recomputed. Actions
remain fixed strata, but their queries are not forced to have equal action
weights. Notebook 3 demonstrates this with small arrays.

\subsection{What the interval means and the rule for the claim}
The 95\% percentile interval describes uncertainty under the chosen resampling
scheme. The 95\% is a nominal repeated-sampling coverage target under its
assumptions, not a 95\% probability that the contextual recipe is superior.
The single confirmatory statement is supported only if its \emph{upper}
bound is below zero. An interval spanning zero is inconclusive for that claim;
it does not prove equivalence or that the true effect is exactly zero.

\begin{tabularx}{\linewidth}{@{}Y Y@{}}
\toprule Observed pattern & Permitted interpretation\\\midrule
Upper bound below zero & Evidence of less contextual degradation under the
declared mixture. Check absolute corrupted accuracy separately.\\\addlinespace
Interval includes zero & Insufficient evidence for the declared smaller-drop
claim; no equivalence conclusion.\\\addlinespace
Lower bound above zero & Evidence points toward more contextual degradation
under the declared mixture.\\\bottomrule
\end{tabularx}

\begin{caution}{The interval's scope is conditional}
It conditions on the trained heads, the three fixed seeds, fixed actions,
gallery anchors, query definitions, and realized synthetic corruptions. It
assumes performance groups are exchangeable within an action for this resampling
purpose. Repeated performances from one person can still be dependent.

It does not quantify uncertainty over arbitrary people, new action classes,
different reference clips, unrestricted training seeds, or real camera failures.
Ten thousand bootstrap draws do not create ten thousand independent models.
\end{caution}

Three seeds expose some training variation but do not justify broad population
claims about random initializations. Every seed and condition is reported.
SupCon, exact-official results, and per-cell comparisons remain descriptive;
there are no additional unplanned superiority tests selected from favorable
rows of the 180-cell table.

\subsection{Sensitivity without changing the retrieval problem}
Leave-one-action-out analysis asks whether one action strongly determines the
reported effect. Remove that action's \emph{queries}, retaining the original
gallery and original ranks. If the gallery item were also removed, the remaining
queries would face fewer alternatives, changing the retrieval problem itself.
The ranges and sign changes are descriptive checks, not new confirmatory tests.

Report effect size, clean accuracy, absolute corrupted accuracy, per-seed and
per-condition results, cost, and failures together. A null or adverse contextual
result is successful research if the valid declared experiment is completed.
A blocked experiment with no final comparison is an honest operational outcome,
but it is not a null scientific result.

\studycheck{Why resample whole performance groups but still weight individual
queries? What does an interval crossing zero establish? Why do the three seeds
not become 10,000 seeds through bootstrapping? Answers: Questions 10--12.}
\source{\src{plan/research-plan.v3.md}{Locked analysis, Section 6};
\src{docs/protocol/protocol-v3.md}{Analysis conventions};
\src{docs/results/README.md}{Reproducible result/report contract}.}
